\section{Conclusion}\label{sec:conclusion}

A national social accounting matrix can be rebuilt from publicly
obtainable sources---in
this application, Statistics South Africa's supply and use tables and
Annual Financial Statistics, the South African Reserve Bank's Quarterly
Bulletin series, and the Labour Market Dynamics and Living Conditions
surveys---with its production and commodity structure exact, its
institutional macro structure exact up to one priced unpublished input and
one discrepancy consistent with an unrecorded adjustment, and its
distributional detail
approximated within disclosed limits and closed by biproportional (RAS)
balancing whose every factor is logged. The exercise yields a balanced,
fully regenerable SAM; a catalogue of divergences between record and
artefact, each first detected by deterministic reconciliation and then
interpreted against the source record; and an open toolkit offered as
reproducible infrastructure for official SAM construction. The same
framework, unchanged, audits Kenya's leading SAM to the boundary its
publication format imposes---verifying structure and macro controls, with
a deviation pattern consistent with the accounts' own statistical
discrepancy being absorbed through consumption and trade---and generates
a Dutch macro accounting matrix from three API calls with no benchmark at
all, every residual attributed, then balances it into a SAM with every
adjustment factor logged. The three cases are one argument: SAM
construction can move from an opaque artisanal exercise to an auditable,
executable statistical-production process, and what separates an exactly
replicable SAM from an unauditable one is not method, geography, or
income level, but whether the underlying tables are published
machine-readably and whether construction choices are recorded where they
act. The idea worth carrying away is therefore the
simple one the evidence keeps confirming: cell-level provenance and
machine-readable construction rules, published alongside every SAM, would
make such audits---and such SAMs---routine. The immediate next steps are
the survey item-level household systems, cross-entropy balancing as an
alternative adjustment principle, and extending the Dutch generation to a
full multi-sector SAM across the ESA-transmitting countries.
